%% kapa: exam=1 round=31 year=2020 period=1 form=A subject=civil_law subject_ko=민법 seq=30 number=30
%% source: Q-Net 감정평가사 기출문제 (cst00309) artlSeq=5207219 file=extracted/1008/5207219/A형/A형/2020년도 제31회 감정평가사 1차 시험_1교시_A형.pdf
%% conversion: kordoc 4.24.51 -> tools/gen1.py (auto); review=proofread-sonnet; flags=none
\begin{question}{30}
甲이 乙 명의의 X토지에 대하여 점유취득시효기간을 완성한 경우에 관한 설명으로 옳지 않은 것을 모두 고른 것은? (다툼이 있으면 판례에 따름)
\begin{framedtext}
ㄱ. 甲이 乙에게 X토지의 소유권이전등기를 청구한 후 乙이 그 토지를 丙에게 처분한 경우, 이는 乙이 자신의 소유권을 행사한 것이므로 乙은 甲에게 불법행위책임을 지지 않는다. \par
ㄴ. 甲이 아직 소유권이전등기를 하지 않고 있던 중, 丙이 취득시효가 완성하기 전에 마친 丙 명의의 가등기에 기하여 소유권이전의 본등기를 한 경우에도 甲은 丙에 대하여 시효취득을 주장할 수 있다. \par
ㄷ. 甲으로부터 X토지의 점유를 승계한 丁은 甲의 취득시효완성의 효과를 주장하여 직접 자기에게 소유권이전등기를 청구하지 못한다. \par
\end{framedtext}
\begin{choices}
  \item ㄴ
  \item ㄷ
  \item ㄱ, ㄴ
  \item ㄱ, ㄷ
  \item ㄴ, ㄷ
\end{choices}
\end{question}
\answer{3}
